\subsection{Algorithmic Realization of Flow-Field Evasion}
\label{sec:algorithm}

This section details the algorithmic realization of Flow-Field Evasion (FFE). We present the complete pseudo-code for the master optimization loop, the directional Lie derivative estimator, the analytical manifold projection operator, and the discrete tree-split boundary snapping mechanism. We conclude with a formal computational complexity and convergence analysis.

\subsubsection{Master Optimization Loop}
Algorithm~\ref{alg:ffe_master} outlines the primary execution loop of Flow-Field Evasion. Given an attack flow $x_0 \in \mathcal{M}_{\text{flow}}$ flagged by the tree ensemble ($F(x_0) \ge 0$), FFE iteratively selects the infinitesimal generator that maximizes the reduction in malicious classification score per unit of physical operational cost.

\begin{algorithm}[htbp]
\caption{Flow-Field Evasion (FFE) Master Optimization Loop}
\label{alg:ffe_master}
\footnotesize
\begin{algorithmic}[1]
\REQUIRE Classifier $F: \mathbb{R}^d \to \mathbb{R}$, initial flow $x_0 \in \mathcal{M}_{\text{flow}}$, feature schema $\mathcal{S}$, step budget $M$, cost vector $\vec{c} = [c_1, c_2, c_3, c_4]^T$, target threshold $\tau = 0.0$
\ENSURE Evasive flow $x^* \in \mathcal{M}_{\text{flow}}$ with $F(x^*) < \tau$, or failure flag
\STATE $x \leftarrow \text{ProjectManifoldInvariants}(x_0, \mathcal{S})$
\STATE $f_{\text{curr}} \leftarrow F(x)$ \COMMENT{Initial model evaluation}
\IF{$f_{\text{curr}} < \tau$}
    \STATE \RETURN $x, \text{SUCCESS}$ \COMMENT{Flow already classified as Benign}
\ENDIF
\STATE $\text{queries} \leftarrow 1$
\FOR{$\text{step} = 1$ \TO $M$}
    \IF{$f_{\text{curr}} < \tau$}
        \STATE \RETURN $x, \text{SUCCESS}$ \COMMENT{Benign classification reached; no snapping needed}
    \ENDIF
    \STATE $\mathcal{D}, q_{\text{used}} \leftarrow \text{EstimateLieDerivatives}(F, x, \mathcal{S})$
    \STATE $\text{queries} \leftarrow \text{queries} + q_{\text{used}}$
    \STATE $k^* \leftarrow \text{null}$; $\text{max\_eff} \leftarrow 0.0$
    \FOR{\textbf{each } $(k, \nabla_k, \vec{v}_k, \Delta s_k, f_k) \in \mathcal{D}$}
        \IF{$\nabla_k < 0$}
            \STATE $\text{eff} \leftarrow (-\nabla_k) / c_k$ \COMMENT{Derivative is negative (score decreases)}
            \IF{$\text{eff} > \text{max\_eff}$}
                \STATE $\text{max\_eff} \leftarrow \text{eff}$; $k^* \leftarrow (\vec{v}_k, \Delta s_k, f_k)$
            \ENDIF
        \ENDIF
    \ENDFOR
    \IF{$k^* \neq \text{null}$}
        \STATE $(\vec{v}, \Delta s, f_{k^*}) \leftarrow k^*$
        \STATE $x \leftarrow \text{ProjectManifoldInvariants}(x + \Delta s \cdot \vec{v}, \mathcal{S})$
        \STATE $f_{\text{curr}} \leftarrow f_{k^*}$ \COMMENT{Score cached from the candidate probe (no extra query)}
    \ELSE
        \STATE \COMMENT{Stagnation: all probed derivatives are non-negative}
        \STATE $x_{\text{snap}} \leftarrow \text{TreeBoundarySnap}(F, x, \mathcal{S})$; $f_{\text{snap}} \leftarrow F(x_{\text{snap}})$
        \IF{$f_{\text{snap}} < f_{\text{curr}}$}
            \STATE $x \leftarrow x_{\text{snap}}$; $f_{\text{curr}} \leftarrow f_{\text{snap}}$
        \ELSE
            \STATE \COMMENT{Fallback Trajectory: Low-and-Slow Temporal Dilation (capped below 120s exporter timeout)}
            \STATE $\vec{v}_t \leftarrow \vec{v}_{\text{time}}(x)$
            \STATE $\Delta t \leftarrow \min(1.1 \times 10^8, \max(50000.0, x_6 \cdot 1.5))$
            \STATE $x \leftarrow \text{ProjectManifoldInvariants}(x + \Delta t \cdot \vec{v}_t, \mathcal{S})$
            \STATE $f_{\text{curr}} \leftarrow F(x)$; $\text{queries} \leftarrow \text{queries} + 1$
        \ENDIF
    \ENDIF
\ENDFOR
\IF{$f_{\text{curr}} < \tau$}
    \STATE \RETURN $x, \text{SUCCESS}$
\ELSE
    \STATE \RETURN $x, \text{FAILURE}$
\ENDIF
\end{algorithmic}
\end{algorithm}

\subsubsection{Adaptive Lie Derivative Numerical Estimator}
Algorithm~\ref{alg:lie_estimator} details the numerical estimation of the directional Lie derivative $\mathcal{L}_{\vec{v}_k} F(x)$ across the non-differentiable step boundaries of the tree ensemble. Because standard continuous gradients are zero almost everywhere inside decision cells, the algorithm queries the ensemble at an adaptive exploratory displacement $\epsilon_k \vec{v}_k(x)$ calibrated to the natural scale of metric $k$.

\begin{algorithm}[htbp]
\caption{Adaptive Directional Lie Derivative Estimator}
\label{alg:lie_estimator}
\footnotesize
\begin{algorithmic}[1]
\REQUIRE Classifier $F$, current flow state $x$, feature schema $\mathcal{S}$
\ENSURE List of candidate evaluation tuples $\mathcal{D} = \{(k, \nabla_k, \vec{v}_k, \Delta s_k, f_k)\}_{k=1}^4$ (with $f_k = F(x_{\text{test}})$ the probed score), total exploratory queries $q$
\STATE $\mathcal{D} \leftarrow \emptyset$
\STATE $q \leftarrow 0$ \COMMENT{Counter for candidate exploratory probes}

\STATE \COMMENT{1. Evaluate Temporal Dilation Generator}
\STATE $\vec{v}_1 \leftarrow \vec{v}_{\text{time}}(x)$
\STATE $\Delta s_1 \leftarrow \max(1000.0, x_6 \cdot 0.50)$ \COMMENT{Adaptive temporal displacement in $\mu$s}
\STATE $x_{\text{test}} \leftarrow \text{ProjectManifoldInvariants}(x + \Delta s_1 \cdot \vec{v}_1, \mathcal{S})$
\STATE $\Delta t_1 \leftarrow \Delta s_1 / 10^6$ \COMMENT{Convert duration step to seconds}
\STATE $f_1 \leftarrow F(x_{\text{test}})$; $\nabla_1 \leftarrow (f_1 - F(x)) / \Delta t_1$ \COMMENT{Derivative per second of delay}
\STATE Append $(1, \nabla_1, \vec{v}_1, \Delta s_1, f_1)$ to $\mathcal{D}$
\STATE $q \leftarrow q + 1$

\STATE \COMMENT{2. Evaluate Payload Padding Generator}
\STATE $\vec{v}_2 \leftarrow \vec{v}_{\text{pad}}(x)$
\STATE $\Delta s_2 \leftarrow \max(50.0, x_4 \cdot 0.25)$ \COMMENT{Adaptive forward byte padding step}
\STATE $x_{\text{test}} \leftarrow \text{ProjectManifoldInvariants}(x + \Delta s_2 \cdot \vec{v}_2, \mathcal{S})$
\STATE $f_2 \leftarrow F(x_{\text{test}})$; $\nabla_2 \leftarrow (f_2 - F(x)) / \Delta s_2$
\STATE Append $(2, \nabla_2, \vec{v}_2, \Delta s_2, f_2)$ to $\mathcal{D}$
\STATE $q \leftarrow q + 1$

\STATE \COMMENT{3. Evaluate Chaff Packet Injection Generator}
\STATE $\vec{v}_3 \leftarrow \vec{v}_{\text{chaff}}(x)$
\STATE $\Delta s_3 \leftarrow \max(1.0, \text{round}(x_2 \cdot 0.50))$ \COMMENT{Adaptive chaff packet injection count}
\STATE $x_{\text{test}} \leftarrow \text{ProjectManifoldInvariants}(x + \Delta s_3 \cdot \vec{v}_3, \mathcal{S})$
\STATE $f_3 \leftarrow F(x_{\text{test}})$; $\nabla_3 \leftarrow (f_3 - F(x)) / \Delta s_3$
\STATE Append $(3, \nabla_3, \vec{v}_3, \Delta s_3, f_3)$ to $\mathcal{D}$
\STATE $q \leftarrow q + 1$

\STATE \COMMENT{4. Evaluate TCP Window Morphing Generator}
\IF{$x[1] = 6$}
    \STATE \COMMENT{TCP flows only; window size is inapplicable to UDP}
    \STATE $\vec{v}_4 \leftarrow \vec{v}_{\text{win}}(x)$
    \STATE $\Delta s_4 \leftarrow |64240.0 - x[7]|$
    \IF{$\Delta s_4 > 0$}
        \STATE $x_{\text{test}} \leftarrow \text{ProjectManifoldInvariants}(x + \Delta s_4 \cdot \vec{v}_4, \mathcal{S})$
        \STATE $f_4 \leftarrow F(x_{\text{test}})$; $\nabla_4 \leftarrow (f_4 - F(x)) / (\Delta s_4 / 65535.0)$ \COMMENT{Normalized window gradient}
        \STATE Append $(4, \nabla_4, \vec{v}_4, \Delta s_4, f_4)$ to $\mathcal{D}$
        \STATE $q \leftarrow q + 1$
    \ENDIF
\ENDIF
\STATE \RETURN $\mathcal{D}, q$
\end{algorithmic}
\end{algorithm}

\subsubsection{Analytical Manifold Projection Operator \texorpdfstring{$\mathcal{P}_{\mathcal{M}}$}{P\_M}}
Algorithm~\ref{alg:manifold_proj} specifies the closed-form projection operator $\mathcal{P}_{\mathcal{M}}: \mathbb{R}^{12} \to \mathcal{M}_{\text{flow}}$. This operator ensures that discrete count boundaries, protocol packet and header length inequalities, and coupled rate equations are strictly satisfied after any continuous displacement.

\begin{algorithm}[htbp]
\caption{Analytical Manifold Projection Operator $\mathcal{P}_{\mathcal{M}}$}
\label{alg:manifold_proj}
\begin{algorithmic}[1]
\REQUIRE Perturbed vector $x \in \mathbb{R}^{12}$, feature schema $\mathcal{S}$, optional baseline flow $x_{\text{orig}}$
\ENSURE Projected flow vector $x_{\text{proj}} \in \mathcal{M}_{\text{flow}}$
\STATE $x_{\text{proj}} \leftarrow \text{copy}(x)$
\IF{$x_{\text{orig}}$ is provided}
    \STATE $x_{\text{proj}}[j] \leftarrow x_{\text{orig}}[j]$ for all $j \in \mathcal{I}_{\text{immut}}$ \COMMENT{Preserve immutable features}
    \STATE $x_{\text{proj}}[j] \leftarrow \max(x_{\text{orig}}[j], x_{\text{proj}}[j])$ for all $j \in \mathcal{I}_{\text{mono}}^+$ \COMMENT{Enforce monotonicity lower bounds}
\ENDIF
\STATE \COMMENT{Enforce discrete bounds and protocol-dependent packet/header sizes}
\STATE $x_{\text{proj}}[2] \leftarrow \max(1.0, \text{round}(x_{\text{proj}}[2]))$ \COMMENT{Integer Forward Packets}
\STATE $x_{\text{proj}}[3] \leftarrow x[3]$ \COMMENT{Immutable Backward Packets Preserved}
\STATE $H_{\text{min}} \leftarrow \mathbf{if}\; x_{\text{proj}}[1] = 6 \;\mathbf{then}\; 20.0 \;\mathbf{else}\; 8.0$ \COMMENT{Min transport header: 20B TCP, 8B UDP}
\STATE $H_{\text{max}} \leftarrow \mathbf{if}\; x_{\text{proj}}[1] = 6 \;\mathbf{then}\; 60.0 \;\mathbf{else}\; 8.0$ \COMMENT{Max transport header: 60B TCP, 8B UDP}
\STATE $x_{\text{proj}}[11] \leftarrow \min(x_{\text{proj}}[2] \times H_{\text{max}}, \max(x_{\text{proj}}[2] \times H_{\text{min}}, \text{round}(x_{\text{proj}}[11])))$ \COMMENT{Transport Header Length}
\STATE $x_{\text{proj}}[4] \leftarrow \min(x_{\text{proj}}[2] \times 1500.0, \max(x_{\text{proj}}[11] + 20.0 \times x_{\text{proj}}[2], \text{round}(x_{\text{proj}}[4])))$ \COMMENT{Forward Bytes $\in [H_{\text{fwd}} + 20 N_{\text{fwd}}, 1500 N_{\text{fwd}}]$}
\STATE $x_{\text{proj}}[5] \leftarrow x[5]$ \COMMENT{Immutable Backward Bytes Preserved}
\STATE $x_{\text{proj}}[6] \leftarrow \min(1.1 \times 10^8, \max(1.0, \text{round}(x_{\text{proj}}[6])))$ \COMMENT{Integer Duration $\ge 1 \mu$s, capped below flow timeout}
\STATE $x_{\text{proj}}[7] \leftarrow \mathbf{if}\; x_{\text{proj}}[1] = 6 \;\mathbf{then}\; \text{clamp}(\text{round}(x_{\text{proj}}[7]), 0.0, 65535.0) \;\mathbf{else}\; 0.0$ \COMMENT{TCP Window $[0, 65535]$, 0 for UDP}

\STATE \COMMENT{Analytically recalculate coupled non-linear flow rates}
\STATE $T_{\text{sec}} \leftarrow x_{\text{proj}}[6] / 10^6$
\STATE $B_{\text{tot}} \leftarrow x_{\text{proj}}[4] + x_{\text{proj}}[5]$
\STATE $N_{\text{tot}} \leftarrow x_{\text{proj}}[2] + x_{\text{proj}}[3]$
\STATE $x_{\text{proj}}[8] \leftarrow B_{\text{tot}} / T_{\text{sec}}$ \COMMENT{Flow Bytes/s}
\STATE $x_{\text{proj}}[9] \leftarrow N_{\text{tot}} / T_{\text{sec}}$ \COMMENT{Flow Packets/s}
\STATE $x_{\text{proj}}[10] \leftarrow x_{\text{proj}}[4] / x_{\text{proj}}[2]$ \COMMENT{Forward Packet Length Mean}
\STATE \RETURN $x_{\text{proj}}$
\end{algorithmic}
\end{algorithm}

\subsubsection{Tree-Split Boundary Snapping (TSBS)}
Algorithm~\ref{alg:tsbs} implements Tree-Split Boundary Snapping. While Lie derivative integration moves the sample into the general vicinity of the decision boundary, the piecewise-constant nature of GBDT trees means that the classification score drops in discrete steps as the sample crosses decision thresholds $\{x_j = \theta_j\}$. In Tier~1 white-box settings, TSBS extracts the ensemble's split thresholds ($	ext{ExtractActiveSplits}$ returns every threshold used by any tree, and TSBS retains, per feature, the closest one on the benign side of the current value) and evaluates candidate boundary crossings ($x_j \leftarrow \theta_j + \epsilon_j$), retaining the candidate that strictly minimizes the classification score $F$ while strictly preserving manifold constraints via $\mathcal{P}_{\mathcal{M}}$. 

Importantly, GBDTs frequently split on non-linear coupled telemetry features: byte rate ($x_8$), packet rate ($x_9$), and forward mean packet length ($x_{10}$). TSBS inverts these coupled constraints in closed form into admissible infinitesimal generator controls:
\begin{enumerate}
    \item \emph{Coupled Rate Inversion:} To cross an active byte rate threshold $x_8 < \theta_8$ or packet rate threshold $x_9 < \theta_9$, TSBS inverts the hyperbolic relation $x_8 = 10^6 (x_4 + x_5) / x_6$ to calculate the required duration dilation:
    \begin{equation}
        x_6^* = \frac{10^6 (x_4 + x_5)}{\theta_8} + \epsilon_6, \quad \text{or} \quad x_6^* = \frac{10^6 (x_2 + x_3)}{\theta_9} + \epsilon_6
    \end{equation}
    with $\epsilon_6 = 10\,\mu$s; since $\theta < x_8$ (resp.\ $x_9$), $x_6^* > x_6$ automatically, and the candidate is discarded if $x_6^*$ exceeds the $1.1 \times 10^8\,\mu$s timeout cap.
    \item \emph{Coupled Packet Length Inversion:} To cross an active packet length threshold $x_{10} < \theta_{10}$ from above, TSBS inverts $x_{10} = (x_4 + N \cdot L_0) / (x_2 + N) < \theta_{10}$, yielding the minimal integer chaff packet count (the smallest integer strictly greater than $(x_4 - \theta_{10} x_2)/(\theta_{10} - L_0)$):
    \begin{equation}
        N_{\text{chaff}} = \left\lfloor \frac{x_4 - \theta_{10} x_2}{\theta_{10} - L_0} \right\rfloor + 1 \quad (\text{for } L_0 < \theta_{10} < x_{10})
    \end{equation}
    \item \emph{Window Advertisement Snapping:} For TCP flows with non-standard window advertisements, TSBS evaluates candidate snaps toward benign OS baselines ($x_7^* = 64,240$) and adjacent split partitions.
\end{enumerate}

\begin{algorithm}[htbp]
\caption{Tree-Split Boundary Snapping (TSBS)}
\label{alg:tsbs}
\footnotesize
\begin{algorithmic}[1]
\REQUIRE Classifier $F$, flow candidate $x \in \mathcal{M}_{\text{flow}}$ with $F(x) \ge \tau$ (still flagged malicious), feature schema $\mathcal{S}$, tolerance $\epsilon_{\text{rel}}$
\ENSURE Snapped flow $x^* \in \mathcal{M}_{\text{flow}}$
\STATE $x^* \leftarrow \text{copy}(x)$; $\mathcal{C} \leftarrow \emptyset$ \COMMENT{Candidate boundary-crossing pool}
\STATE $\text{active\_splits} \leftarrow \text{ExtractActiveSplits}(F, x^*)$
\STATE \COMMENT{1. Monotonic features: closest split strictly above $x^*[j]$}
\FOR{\textbf{each } $j \in \mathcal{I}_{\text{mono}}^+$}
    \STATE $\theta_j^* \leftarrow \min\{\theta_j \mid (j, \theta_j) \in \text{active\_splits} \;\land\; \theta_j > x^*[j]\}$
    \IF{$\theta_j^*$ exists}
        \IF{$j = 11$ \AND $x^*[1] = 17$}
            \STATE \textbf{continue} \COMMENT{UDP transport header immutable at 8B}
        \ENDIF
        \STATE $\epsilon_j \leftarrow \mathbf{if}\; j \in \{2, 4, 11\} \;\mathbf{then}\; 1.0 \;\mathbf{else}\; \text{round}(\epsilon_{\text{rel}} \cdot \max(1.0, |\theta_j^*|))$
        \STATE $x_{\text{c}} \leftarrow \text{copy}(x^*)$; $x_{\text{c}}[j] \leftarrow \theta_j^* + \epsilon_j$
        \IF{$j = 11$}
            \STATE $x_{\text{c}}[11] \leftarrow \min(60.0 \cdot x_{\text{c}}[2], x_{\text{c}}[11])$
            \STATE $x_{\text{c}}[4] \leftarrow \min(1500.0 \cdot x_{\text{c}}[2], \max(x_{\text{c}}[4], x_{\text{c}}[11] + 20.0 \cdot x_{\text{c}}[2]))$
        \ENDIF
        \STATE $\mathcal{C} \leftarrow \mathcal{C} \cup \{x_{\text{c}}\}$
    \ENDIF
\ENDFOR
\STATE \COMMENT{2. Invert coupled rate splits ($x_8, x_9$) into duration dilation}
\STATE $\theta_8^* \leftarrow \max\{\theta_8 \mid (8, \theta_8) \in \text{active\_splits} \;\land\; 0 < \theta_8 < x^*[8]\}$
\IF{$\theta_8^*$ exists}
    \STATE $d \leftarrow (10^6 \cdot (x^*[4] + x^*[5]) / \theta_8^*) + 10.0$
    \IF{$d \le 1.1 \times 10^8$}
        \STATE $x_{\text{c}} \leftarrow \text{copy}(x^*)$; $x_{\text{c}}[6] \leftarrow d$; $\mathcal{C} \leftarrow \mathcal{C} \cup \{x_{\text{c}}\}$
    \ENDIF
\ENDIF
\STATE $\theta_9^* \leftarrow \max\{\theta_9 \mid (9, \theta_9) \in \text{active\_splits} \;\land\; 0 < \theta_9 < x^*[9]\}$
\IF{$\theta_9^*$ exists}
    \STATE $d \leftarrow (10^6 \cdot (x^*[2] + x^*[3]) / \theta_9^*) + 10.0$
    \IF{$d \le 1.1 \times 10^8$}
        \STATE $x_{\text{c}} \leftarrow \text{copy}(x^*)$; $x_{\text{c}}[6] \leftarrow d$; $\mathcal{C} \leftarrow \mathcal{C} \cup \{x_{\text{c}}\}$
    \ENDIF
\ENDIF
\STATE \COMMENT{3. Invert coupled packet length splits ($x_{10}$) into chaff packet injection}
\STATE $H_0 \leftarrow \mathbf{if}\; (x^*[1] = 6 \;\land\; x^*[11]/x^*[2] \ge 26.0) \;\mathbf{then}\; 32.0 \;\mathbf{else}\; (\mathbf{if}\; x^*[1]=6 \;\mathbf{then}\; 20.0 \;\mathbf{else}\; 8.0)$
\STATE $L_0 \leftarrow 20.0 + H_0$ \COMMENT{RFC 7323 options compliance: 52B/40B for TCP, 28B for UDP}
\STATE $\theta_{10}^* \leftarrow \max\{\theta_{10} \mid (10, \theta_{10}) \in \text{active\_splits} \;\land\; L_0 < \theta_{10} < x^*[10]\}$
\IF{$\theta_{10}^*$ exists}
    \STATE $N \leftarrow \lfloor (x^*[4] - \theta_{10}^* x^*[2]) / (\theta_{10}^* - L_0) \rfloor + 1$
    \IF{$N \le 1000$}
    \STATE $x_{\text{c}} \leftarrow \text{copy}(x^*)$; $x_{\text{c}}[2] \leftarrow x^*[2] + N$; $x_{\text{c}}[4] \leftarrow x^*[4] + N \cdot L_0$; $x_{\text{c}}[11] \leftarrow x^*[11] + N \cdot H_0$
    \STATE $\mathcal{C} \leftarrow \mathcal{C} \cup \{x_{\text{c}}\}$ \COMMENT{Chaff budget cap: $N \le 1000$ packets}
    \ENDIF
\ENDIF
\STATE \COMMENT{4. TCP Window morphing to benign OS baseline}
\IF{$x^*[1] = 6$ \AND $x^*[7] \neq 64240.0$}
    \STATE $x_{\text{c}} \leftarrow \text{copy}(x^*)$; $x_{\text{c}}[7] \leftarrow 64240.0$; $\mathcal{C} \leftarrow \mathcal{C} \cup \{x_{\text{c}}\}$
\ENDIF
\STATE \COMMENT{Evaluate projected candidates and retain best partition crossing}
\STATE $x_{\text{best}} \leftarrow x^*$; $s_{\text{best}} \leftarrow F(x^*)$
\FOR{\textbf{each } $x_{\text{c}} \in \mathcal{C}$}
    \STATE $x_{\text{proj}} \leftarrow \text{ProjectManifoldInvariants}(x_{\text{c}}, \mathcal{S})$
    \IF{$F(x_{\text{proj}}) < s_{\text{best}}$}
        \STATE $s_{\text{best}} \leftarrow F(x_{\text{proj}})$; $x_{\text{best}} \leftarrow x_{\text{proj}}$
        \IF{$s_{\text{best}} < 0$}
            \STATE \textbf{break} \COMMENT{Early exit once benign classification is achieved}
        \ENDIF
    \ENDIF
\ENDFOR
\STATE \RETURN $x_{\text{best}}$
\end{algorithmic}
\end{algorithm}

\subsubsection{Computational Complexity \& Comparison to MILP}
Table~\ref{tab:complexity_comparison} compares the asymptotic time and space complexity of Flow-Field Evasion against exact MILP solvers \cite{kantchelian2016evasion} and heuristic black-box optimizers (MOEVA \cite{ghamizi2020moeva}, CAA \cite{simonetto2024constrained}, Traffic Manipulator \cite{han2021traffic}).

\begin{table*}[htbp]
\caption{Asymptotic Complexity Comparison of Evasion Algorithms for Tree Ensembles}\label{tab:complexity_comparison}%
\centering
\resizebox{\textwidth}{!}{%
\begin{tabular}{@{}lccc@{}}
\toprule
\textbf{Algorithm} & \textbf{Time Complexity} & \textbf{Query Complexity} & \textbf{Coupled Invariants} \\
\midrule
\textbf{Kantchelian MILP} \cite{kantchelian2016evasion} & $2^{\mathcal{O}(D \cdot T)}$ (Worst Case) & $0$ (White-Box) & Unsupported (Linear Only) \\
\textbf{MOEVA / CAA} \cite{simonetto2024constrained, ghamizi2020moeva}   & $\mathcal{O}(G \cdot P \cdot T \cdot D)$   & $\mathcal{O}(10^3 - 10^4)$ & Approximate Penalty \\
\textbf{Traffic Manipulator} \cite{han2021traffic}    & $\mathcal{O}(I \cdot S \cdot T \cdot D)$   & $\mathcal{O}(10^3 - 10^4)$ & Packet Re-evaluation \\
\midrule
\textbf{Flow-Field Evasion (Ours)}                     & $\mathcal{O}(M \cdot K \cdot T \cdot D)$ & $\le 5M_{\max} + 1$ external (mean $12.9$ incl.\ TSBS) & \textbf{Exact Closed-Form} \\
\bottomrule
\end{tabular}%
}
\end{table*}

Where:
\begin{itemize}
    \item $T$ is the number of trees in the ensemble (typically $50 - 100$).
    \item $D$ is the maximum depth of each tree (typically $6 - 10$).
    \item $M_{\max}$ is the maximum FFE iteration budget ($M_{\max} = 20$).
    \item $K$ is the number of infinitesimal Lie generators ($K = 4$).
    \item $G, P$ are the generations and population size in MOEVA ($G = 25, P = 50$, totaling $\le 1,250$ queries).
    \item $I, S$ are the iterations and swarm size in PSO ($I = 37, S = 50$, totaling $\le 1,850$ queries).
\end{itemize}

\paragraph{Query Accounting Details.}
In the worst case, executing all $M_{\max} = 20$ iterations without early termination requires $1$ baseline query plus up to $K = 4$ exploratory probes per iteration, giving an upper bound of $Q_{\max} \le 4 M_{\max} + 1 = 81$ model queries (or $\le 101$ queries if fallback dilation evaluations are triggered on stagnation). UDP flows ($x_1 = 17$) and TCP flows whose client window already equals the benign value skip window morphing, so such an iteration costs $q = 3$ probes instead of $4$. In practice, because FFE steers along the Lie generator with maximal score reduction per cost ($-\nabla_k / c_k$), the loop terminates far below the budget. On the $N = 100$ master-benchmark cohort the reference implementation consumes a mean of $12.9$ model queries per flow, of which $10.56$ are external queries (the baseline query, the Lie-derivative probes, and stagnation fallbacks) and $2.38$ are candidate evaluations inside the TSBS routine (Algorithm~\ref{alg:tsbs}).

In Tier~1 white-box settings, TSBS extracts candidate thresholds directly from the adversary's local copy of the ensemble ($\text{ExtractActiveSplits}$), so its candidate evaluations could be executed offline against that copy. The reported figure of $12.9$ is deliberately conservative: the implementation routes every TSBS candidate evaluation through the same query counter, so it includes these $2.38$ evaluations. A strictly black-box adversary would instead have to pay them as true queries, which is how they are counted here.

Because $K \le 4$ and $M \le 20$ are small constants, FFE requires only a linear number of simple tree evaluations $\mathcal{O}(T \cdot D)$. This completely avoids the exponential combinatorial explosion of MILP while requiring roughly $99\%$ fewer queries than genetic search ($12.9$ vs.\ $\le 1{,}250$) and PSO ($12.9$ vs.\ $\le 1{,}850$).
